\section{最初的推论}

\begin{metatheorem}{假言三段论}{}
	设$\boldsymbol{A},\boldsymbol{B},\boldsymbol{C}$是理论$\mathcal{T}$中的公式.若$\boldsymbol{A}\implies\boldsymbol{B}$和$\boldsymbol{B}\implies\boldsymbol{C}$均是$\mathcal{T}$的定理,则$\boldsymbol{A}\implies\boldsymbol{C}$是$\mathcal{T}$的定理.
\end{metatheorem}

\begin{metatheorem}{右析取引入律}{}
	设$\boldsymbol{A}$和$\boldsymbol{B}$是理论$\mathcal{T}$中的公式,则$\boldsymbol{B}\implies\left(\boldsymbol{A}\vee\boldsymbol{B}\right)$是$\mathcal{T}$的定理.
\end{metatheorem}

\begin{metatheorem}{同一律}{}
	设$\boldsymbol{A}$是理论$\mathcal{T}$中的公式,则$\boldsymbol{A}\implies\boldsymbol{A}$是$\mathcal{T}$的定理.
\end{metatheorem}

\begin{metatheorem}{前件弱化律}{}
	设$\boldsymbol{A}$是理论$\mathcal{T}$中的公式,$\boldsymbol{B}$是$\mathcal{T}$中的一个定理,则$\boldsymbol{A}\implies\boldsymbol{B}$是$\mathcal{T}$的定理.
\end{metatheorem}

\begin{metatheorem}{排中律}{}
	设$\boldsymbol{A}$是理论$\mathcal{T}$中的公式,则$\boldsymbol{A}\vee\neg\boldsymbol{A}$是$\mathcal{T}$的定理.
\end{metatheorem}

\begin{metatheorem}{双重否定引入律}{}
	设$\boldsymbol{A}$是理论$\mathcal{T}$中的公式,则$\boldsymbol{A}\implies\neg\neg\boldsymbol{A}$是$\mathcal{T}$的定理.
\end{metatheorem}

\begin{metatheorem}{逆否命题定理}{}
	设$\boldsymbol{A}$和$\boldsymbol{B}$是理论$\mathcal{T}$中的两个公式,则$\left(\boldsymbol{A}\implies\boldsymbol{B}\right)\implies\left(\neg\boldsymbol{B}\implies\neg\boldsymbol{A}\right)$是$\mathcal{T}$中的定理.
\end{metatheorem}

\begin{metatheorem}{蕴含的传递性}{}
	设$\boldsymbol{A},\boldsymbol{B},\boldsymbol{C}$是理论$\mathcal{T}$中的公式.若$\boldsymbol{A}\implies\boldsymbol{B}$是$\mathcal{T}$中的定理,则$\left(\boldsymbol{B}\implies\boldsymbol{C}\right)\implies\left(\boldsymbol{A}\implies\boldsymbol{C}\right)$是$\mathcal{T}$的定理.
\end{metatheorem}

自此以后,在书写形式证明时,我们将默认并自由地使用肯定前件律和蕴含的传递性进行推理,(通常)不再显式地写出其引用.